% GENERATED FILE. Edit canonical lessons, then run npm run generate:books.
% canonical-source-hash: fnv1a64:ad5a42fac355be29
% canonical-lessons: MW-C153-hu-hai, MW-C153-aap-ho, MW-C153-unt, MW-W45-uu, MW-R153-am-is-are

\chapter{Am, Is, Are}
\label{ch:mw-a2-am-is-are}
\hlchaptermodality{153}

\begin{chapteropening}
\textbf{By the end of this chapter:} \emph{I can say am, is and are with \mw{हूं}, \mw{आप}, \mw{वो} and any noun, and write the letter \mw{ऊ} from \mw{ऊंट}.}

Complete MW-R153-am-is-are: say am, is and are for each person, name the camel, and write \mw{ऊ}.
\end{chapteropening}

\section[\texorpdfstring{\mw{हूं} · \mw{है} (hū̃ · hai)}{हूं · है (hū̃ · hai)}]{\mw{हूं} and \mw{है} — am and is}

\label{lesson:MW-C153-hu-hai}



\begin{quote}
Say \emph{I am fine}, then \emph{he is ill}. You have said \emph{am} and \emph{is} since the greeting; now they get a name.
\end{quote}

\begin{grammarlens}[title={Grammar Lens: \mw{हूं} … \mw{हूं}}]
In \textbf{\mw{हूं} \mw{ठीक} \mw{हूं।}} the first \textbf{\mw{हूं}} is \emph{I} and the last one is \emph{am}: the same word, at both ends.

\begin{quote}
\textbf{\mw{हूं} \mw{अठै} \mw{हूं।}}
\end{quote}

\begin{quote}
\emph{hū̃ aṭhai hū̃.} — I am here.
\end{quote}
\end{grammarlens}

\begin{grammarlens}[title={Grammar Lens: \mw{है}}]
Anyone or anything you talk about takes \textbf{\mw{है}}, \emph{is}: \textbf{\mw{वो}}, \textbf{\mw{राम}}, the tea.

\begin{quote}
\textbf{\mw{राम} \mw{बठै} \mw{है।}}
\end{quote}

\begin{quote}
\emph{rām baṭhai hai.} — Ram is there.
\end{quote}

\begin{quote}
\textbf{\mw{चाय} \mw{गरम} \mw{है।}}
\end{quote}

\begin{quote}
\emph{chāy garam hai.} — The tea is hot.
\end{quote}
\end{grammarlens}

\subsection*{Your turn}
Say these aloud:

\begin{itemize}
\raggedright
  \item \textbf{\mw{हूं} \mw{अठै} \mw{हूं} \mw{पण} \mw{राम} \mw{बठै} \mw{है।}} — \emph{hū̃ aṭhai hū̃ paṇ rām baṭhai hai.} — I am here, but Ram is there.
  \item \textbf{\mw{हूं} \mw{राजी} \mw{हूं} \mw{क्यूंकै} \mw{चाय} \mw{गरम} \mw{है।}} — \emph{hū̃ rājī hū̃ kyū̃kai chāy garam hai.} — I am happy because the tea is hot.
  \item \textbf{\mw{वो} \mw{उदास} \mw{है।}} — \emph{vo udās hai.} — He is sad.
\end{itemize}

\begin{tcolorbox}[breakable,colback=teal!4,colframe=teal!35!black,title={Before you move on}]
\emph{I am} ends on? (\textbf{\mw{हूं}}.) \emph{He is}? (\textbf{\mw{है}}.)
\end{tcolorbox}
\section[\texorpdfstring{\mw{आप} … \mw{हो} (āp … ho)}{आप … हो (āp … ho)}]{\mw{आप} … \mw{हो} — you are}

\label{lesson:MW-C153-aap-ho}



\begin{quote}
Ask the polite question: \textbf{\mw{आप} \mw{कैसो} \mw{हो}?} Then say \emph{I am here} and \emph{Ram is there}.
\end{quote}

\begin{grammarlens}[title={Grammar Lens: \mw{हो} after \mw{आप}}]
\textbf{\mw{आप} \mw{कैसो} \mw{हो}?} ends on \textbf{\mw{हो}}. After \textbf{\mw{आप}}, \emph{are} is \textbf{\mw{हो}}:

\begin{quote}
\textbf{\mw{आप} \mw{अठै} \mw{हो।}}
\end{quote}

\begin{quote}
\emph{āp aṭhai ho.} — You are here.
\end{quote}

\textbf{\mw{हो}} answers to \textbf{\mw{आप}} the way \textbf{\mw{हूं}} answers to \textbf{\mw{हूं}}, and \textbf{\mw{है}} to \textbf{\mw{वो}}.
\end{grammarlens}

\subsection*{Your turn}
Say these aloud:

\begin{itemize}
\raggedright
  \item \textbf{\mw{आप} \mw{बठै} \mw{हो} \mw{पण} \mw{हूं} \mw{अठै} \mw{हूं।}} — \emph{āp baṭhai ho paṇ hū̃ aṭhai hū̃.} — You are there, but I am here.
  \item \textbf{\mw{आप} \mw{राजी} \mw{हो।}} — \emph{āp rājī ho.} — You are happy.
  \item \textbf{\mw{चाय} \mw{अर} \mw{पाणी}?} — \emph{chāy ar pāṇī?} — Tea or water?
  \item \textbf{\mw{पाणी।}} — \emph{pāṇī.} — Water.
\end{itemize}

\begin{tcolorbox}[breakable,colback=teal!4,colframe=teal!35!black,title={Before you move on}]
Which \emph{are} goes with \textbf{\mw{आप}}? (\textbf{\mw{हो}}.) And with \textbf{\mw{वो}}? (\textbf{\mw{है}}.)
\end{tcolorbox}
\section[\texorpdfstring{\mw{ऊंट} (ū̃ṭ)}{ऊंट (ū̃ṭ)}]{\mw{ऊंट} (ū̃ṭ) — a camel}

\label{lesson:MW-C153-unt}



\begin{quote}
Before the new one: say \emph{a cow}, \emph{a horse} and \emph{an elephant}. (\textbf{\mw{गाय}}, \textbf{\mw{घोड़ो}}, \textbf{\mw{हाथी}}.)
\end{quote}

\subsection*{You'll want to know: \mw{ऊंट}}
\textbf{\mw{ऊंट}} — \emph{ū̃ṭ} — \textquotedblleft{}a camel\textquotedblright{}.

\begin{quote}
\textbf{\mw{वो} \mw{ऊंट} \mw{है।}}
\end{quote}

\begin{quote}
\emph{vo ū̃ṭ hai.} — That is a camel.
\end{quote}

It opens on a letter you have not written yet: the long \emph{ū}, standing alone. The next lesson puts it in your hand.

\begin{cousinweb}
\textbf{\mw{ऊंट}} (\emph{ū̃ṭ}) is inherited from Sanskrit \emph{uṣṭra}, which named a buffalo as well as a camel. CDIAL lists Old Marwari \emph{uṭa} in the same entry (CDIAL 2387).
\end{cousinweb}

\begin{grammarlens}[title={What you've built}]
\textbf{\mw{ऊंट}} is masculine, and takes \textbf{\mw{है}} like anything you talk about: \textbf{\mw{ऊंट} \mw{अठै} \mw{है।}} (\emph{ū̃ṭ aṭhai hai.}), the camel is here.
\end{grammarlens}

\subsection*{Your turn}
\begin{itemize}
\raggedright
  \item \emph{Say it:} \emph{ū̃ṭ}
  \item \emph{Say it:} \textbf{\mw{वो} \mw{ऊंट} \mw{है।}} — \emph{vo ū̃ṭ hai.} — That is a camel.
  \item \emph{Recall:} say \emph{a cow}, then \emph{ū̃ṭ} again
\end{itemize}

\begin{tcolorbox}[breakable,colback=teal!4,colframe=teal!35!black,title={Before you move on}]
What does \textbf{\mw{ऊंट}} mean? (\textquotedblleft{}A camel\textquotedblright{}.) Say it once more.
\end{tcolorbox}
\section[\texorpdfstring{\mw{ऊ} (ū)}{ऊ (ū)}]{\mw{ऊ} (ū) — the letter that opens \mw{ऊंट}}

\label{lesson:MW-W45-uu}



\begin{quote}
Before the new one: write \textbf{\mw{उ}} once, and name it. Then say \textbf{\mw{हूं}} (\emph{hū̃}): its long \emph{ū} is the sign \textbf{\mw{ू}}.

Now say \textbf{\mw{ऊंट}} (\emph{ū̃ṭ}), \textquotedblleft{}a camel\textquotedblright{}. It opens on the long vowel standing alone, and this lesson puts that letter in your hand.
\end{quote}

\subsection*{Script you'll notice: \mw{ऊ}}
\textbf{\mw{ऊ}} — \emph{ū}, the long \emph{u}.

It opens \textbf{\mw{ऊंट}}. It is the letter \textbf{\mw{उ}} with one more part: a separate loop on its right. \textbf{\mw{उ}} and \textbf{\mw{ऊ}} pair the way the signs \textbf{\mw{ु}} and \textbf{\mw{ू}} do, short and long.

\subsection*{Writing: \mw{ऊ}}
\hlblockfigure{\detokenize{figures/MW-W45-uu-filmstrip.pdf}}{How \mw{ऊ} is written, stroke by stroke}

\begin{itemize}
\raggedright
  \item \textbf{1.} curve down and left around the upper bowl, then sweep back through the waist and around the lower loop without lifting
  \item \textbf{2.} lift at the waist and sweep the right-hand loop up, around, and down-left
  \item \textbf{3.} lift and draw the top shirorekhā left-to-right
\end{itemize}

\textbf{Pen lifts: 2.}

\begin{quote}
verified modern printed teaching form; everyday handwriting may simplify the curves
\end{quote}

\begin{quote}
Stroke order is one attested teaching order; hands and teachers vary. Source: Saurmandal, ‘Devanagari \mw{ऊ} stroke order.svg’, panels 1–3, Wikimedia Commons, 5 August 2023.
\end{quote}

\subsection*{Your turn}
\begin{itemize}
\raggedright
  \item \emph{Look:} at \textbf{\mw{ऊंट}} and point at \mw{ऊ} inside it
  \item \emph{Trace it:} \mw{ऊ} three times, saying \emph{ū} as you finish each one
  \item \emph{Write it:} \mw{उ}, then \mw{ऊ}, from memory
  \item \emph{Read it:} \textbf{\mw{ऊंट}} aloud once more
\end{itemize}

\begin{tcolorbox}[breakable,colback=teal!4,colframe=teal!35!black,title={Before you move on}]
Which letter is this — \mw{ऊ}? (\textbf{ū}, the long \emph{u}.) Which word did it come from? (\textbf{\mw{ऊंट}}, \emph{ū̃ṭ}, \textquotedblleft{}a camel\textquotedblright{}.)
\end{tcolorbox}
\section[(dialogue)]{Am, is, are — a review}

\label{lesson:MW-R153-am-is-are}



\begin{quote}
Say \emph{I am here}.
\end{quote}

\subsection*{Your turn}
The three forms side by side, now that each one is taught:

\noindent
\begin{tabularx}{\linewidth}{@{}>{\raggedright\arraybackslash}X>{\raggedright\arraybackslash}X@{}}
\toprule
\textbf{who} & \textbf{to be} \\
\midrule
\textbf{\mw{हूं}} & \textbf{\mw{हूं}} \\
\textbf{\mw{आप}} & \textbf{\mw{हो}} \\
\textbf{\mw{वो}}, \textbf{\mw{राम}}, \textbf{\mw{ऊंट}} & \textbf{\mw{है}} \\
\bottomrule
\end{tabularx}

\begin{itemize}
\raggedright
  \item \emph{Say it:} \emph{I am here}; \emph{you are there}; \emph{Ram is there}
  \item \emph{Say it:} \emph{that is a camel}; \emph{the camel is here}
  \item \emph{Say it:} \emph{I am happy because the tea is hot}
  \item \emph{Say it:} \emph{you are there, but I am here}
  \item \emph{Say it:} \emph{tea or water?}, and answer \emph{tea}
  \item \emph{Read it:} \textbf{\mw{ऊंट}}, and name its first letter
\end{itemize}

\begin{tcolorbox}[breakable,colback=teal!4,colframe=teal!35!black,title={Before you move on}]
\emph{You are}? (\textbf{\mw{आप} … \mw{हो}}.) \emph{The camel is}? (\textbf{\mw{ऊंट} … \mw{है}}.) \emph{I am}? (\textbf{\mw{हूं} … \mw{हूं}}.)
\end{tcolorbox}
